\documentclass[12pt, a4paper]{article}

% Paket-Importe
\usepackage[utf8]{inputenc}
\usepackage[ngerman]{babel}
\usepackage{amsmath, amsfonts, amssymb}
\usepackage{graphicx}
\usepackage{hyperref}

\title{Neovim LaTeX-Integrationstest}
\author{Test Benutzer}
\date{\today}

\begin{document}

\maketitle

\section{Einführung}
\label{sec:einfuehrung}
Dies ist ein Testdokument für die Neovim-Integration. Wir testen hier die
Syntag-Hervorhebung, das Kompilieren sowie die Vorwärts- und Rückwärtssuche.

\section{Mathematik-Test}
\label{sec:mathe}
Ein einfacher In-Line-Test: $E = mc^2$ oder $a^2 + b^2 = c^2$.

Hier ist eine abgesetzte Gleichung mit Nummerierung für die LSP-Referenzierung:
\begin{equation}
    \label{eq:einstein}
    \int_{-\infty}^{\infty} e^{-x^2} \, dx = \sqrt{\pi}
\end{equation}

\section{Referenzen \& Navigations-Test}
\label{sec:referenzen}
In Abschnitt \ref{sec:einfuehrung} haben wir begonnen. Wie in Gleichung
\ref{eq:einstein} zu sehen ist, funktioniert das Mathe-Set hervorragend.

\section{LSP-Diagnose-Test (Syntax-Fehler)}
% Der folgende Befehl löst bei aktiviertem ChkTeX / TexLab eine Linter-Warnung aus,
% da vor dem Punkt ein Leerzeichen steht :
Das ist ein Test für die Fehlersuche .

\end{document}
